% ================================================================
%  conteudo/semana01.tex
% ================================================================

\section{Semana 1 --- 16 de março de 2026 a 20 de março de 2026}

\begin{table}[H]
\centering
\caption{Registro de atividades da Semana 1.}
\begin{tabular}{llll}
\toprule
\textbf{Data} & \textbf{Dia} & \textbf{Horas} & \textbf{Descrição} \\
\midrule
16/03/2026 & Segunda-feira & 2,0 h & Atividades no CIA \\
17/03/2026 & Terça-feira & 2,0 h & Atividades no CIA \\
18/03/2026 & Quarta-feira & 4,0 h & Atividades no CIA \\
19/03/2026 & Quinta-feira & 4,0 h & Atividades no CIA \\
20/03/2026 & Sexta-feira & 5,5 h & Atividades no CIA \\
\midrule
\multicolumn{2}{l}{\textbf{Total da semana}} & \textbf{17,5 h} & \\
\multicolumn{2}{l}{\textbf{Total acumulado}} & \textbf{17,5 h} & \\
\bottomrule
\end{tabular}
\end{table}

\subsubsection*{Atividades realizadas}

Durante a primeira semana de estágio, foram realizadas atividades de
\textbf{análise e preparo de materiais}, com ambientação nos procedimentos
operacionais do CIA/FURG. As principais atividades desenvolvidas foram:

\begin{itemize}
    \item Análise e preparo de materiais para as seguintes técnicas analíticas:
    \begin{itemize}
        \item \textbf{IRMS} --- Espectrometria de Massas de Razão Isotópica;
        \item \textbf{TGA} --- Termogravimetria;
        \item \textbf{RMN} --- Ressonância Magnética Nuclear;
        \item \textbf{BET} --- Adsorção de nitrogênio para determinação de
              superfície e porosidade;
        \item \textbf{DSC} --- Calorimetria Exploratória de Varredura.
    \end{itemize}
    \item Operação dos equipamentos: todos os sistemas são automatizados ---
          o operador insere o material, configura os parâmetros no \textit{software}
          e aguarda a geração automática do relatório analítico;
    \item \textbf{Manutenção e reabastecimento dos gases} utilizados nos
          equipamentos (gases de arraste, gases criogênicos e gases de purga).
\end{itemize}

\subsubsection*{Observações}

A semana foi dedicada à ambientação no ambiente laboratorial do CIA e ao
aprendizado dos fluxos operacionais de cada equipamento. Destaca-se o
caráter altamente automatizado das análises, em que o papel do operador
concentra-se no preparo adequado das amostras, na configuração correta
dos parâmetros e na manutenção preventiva dos equipamentos, incluindo o
controle e reabastecimento dos gases analíticos.
